%  LaTeX support: latex@mdpi.com 
%  For support, please attach all files needed for compiling as well as the log file, and specify your operating system, LaTeX version, and LaTeX editor.

%=================================================================
\documentclass[technologies,article,submit,pdftex,moreauthors]{Definitions/mdpi} 

%=================================================================

% MDPI internal commands
\firstpage{1} 
\makeatletter 
\setcounter{page}{\@firstpage} 
\makeatother
\pubvolume{1}
\issuenum{1}
\articlenumber{0}
\pubyear{2023}
\copyrightyear{2023}
%\externaleditor{Academic Editor: Firstname Lastname}
\datereceived{XX January 2023} 
%\daterevised{} % Only for the journal Acoustics
\dateaccepted{} 
\datepublished{} 
\hreflink{https://doi.org/} % If needed use \linebreak
%\doinum{}

%=================================================================
% Add packages and commands here. The following packages are loaded in our class file: fontenc, inputenc, calc, indentfirst, fancyhdr, graphicx, epstopdf, lastpage, ifthen, lineno, float, amsmath, setspace, enumitem, mathpazo, booktabs, titlesec, etoolbox, tabto, xcolor, soul, multirow, microtype, tikz, totcount, changepage, attrib, upgreek, cleveref, amsthm, hyphenat, natbib, hyperref, footmisc, url, geometry, newfloat, caption

\graphicspath{{figures/}} % path to folder with figures
\usepackage{subcaption}
\usepackage{listings}
\usepackage{gensymb} % for \textdegree
\usepackage{tabularx}

%=================================================================
% Full title of the paper (Capitalized)
\Title{GDAL and PROJ libraries integrated with GRASS GIS for terrain modelling and projecting the georeferenced raster image}

% MDPI internal command: Title for citation in the left column
\TitleCitation{GDAL and PROJ libraries integrated with GRASS GIS for terrain modelling and projecting the georeferenced raster image}

% Author Orchid ID: enter ID or remove command
\newcommand{\orcidauthorA}{0000-0002-5759-1089} % Polina Add \orcidA{} behind the author's name
\newcommand{\orcidauthorB}{0000-0002-6461-1551} % Olivier Add \orcidB{} behind the author's name

% Authors, for the paper (add full first names)
\Author{Polina Lemenkova $^{1}$*\orcidA{} and Olivier Debeir $^{1}$\orcidB{}}

%\longauthorlist{yes}

% MDPI internal command: Authors, for metadata in PDF
\AuthorNames{Polina Lemenkova and Olivier Debeir}

% MDPI internal command: Authors, for citation in the left column
\AuthorCitation{Lemenkova, P.; Debeir, O.}
% If this is a Chicago style journal: Lastname, Firstname, Firstname Lastname, and Firstname Lastname.

% Affiliations / Addresses (Add [1] after \address if there is only one affiliation.)
\address{%
$^{1}$ \quad Laboratory of Image Synthesis and Analysis (LISA), École polytechnique de Bruxelles (Brussels Faculty of Engineering), Université Libre de Bruxelles (ULB). Building L, Campus du Solbosch, ULB -- LISA CP165/57, Avenue Franklin Roosevelt 50, Brussels 1000, Belgium.
}

% Contact information of the corresponding author
\corres{Correspondence: polina.lemenkova@ulb.be; Tel.: +32 471 86 04 59 (P.L.)}

% Current address and/or shared authorship
%\firstnote{Current address: Affiliation 3.} 

% Abstract (Do not insert blank lines, i.e. \\) 
\abstract{Libraries with prewritten codes optimize the workflow in cartography and reduce labour intensive data processing by iteratively applying scripts to implementing mapping tasks. Most existing GIS approaches are based on traditional software with a graphical user's interface which significantly limits their performance. Although plugins are proposed to improve the functionality of many GIS, they are usually ad hoc in finding specific mapping solutions, e.g., cartographic projections and data conversion. We address this limitation by applying the principled approach of GDAL, PROJ and GRASS GIS for geospatial data processing and morphometric analysis. This research presents topographic analysis of the dataset using scripting methods which include several tools: 1) GDAL, a translator library for raster and vector geospatial data formats used for converting ETOPO1 GeoTIFF in XY Cartesian coordinates into WGS84 by the ‘gdalwarp’ utility; 2) PROJ projection transformation library used for converting ETOPO1 WGS84 grid to cartographic projections (Cassini-Soldner equirectangular; Equal Area Cylindrical; 2-point Equidistant Azimuthal; Oblique Mercator); 3) GRASS GIS by sequential use of the following modules: r.info, d.mon, d.rast, r.colors, d.rast.leg, d.legend, d.northarrow, d.grid, d.text, g.region, r.contour. The depth frequency was analysed by the module 'd.histogram'. The proposed approach provided a systematic way for morphometric measuring of topographic data and combining the advantages of the GDAL, PROJ and GRASS GIS tools that include the informativeness, effectiveness and representativeness in spatial data processing. The morphometric analysis included the computed slope, aspect, profile and tangential curvature of the study area. The data analysis revealed the distribution pattern in topographic data: 24\% of data with elevations below 400 m; 13\% of data with depths -5,000 to -6,000 m; 4\% of depths have values -3,000 to -4,000 m, the least frequent data (-6,000 to 7,000 m) <1\%, 2\% of depths have values -2,000 to 3,000 m in the basin, while other values are distributed proportionally. Further, by incorporating the generic coordinate transformation software library PROJ, we transformed raster grid into various cartographic projections to demonstrate distortions in shape and area. Scripting techniques of GRASS GIS are demonstrated for applications in topographic modelling and raster data processing. The GRASS GIS shows the effectiveness for mapping and visualization, compatibility with libraries (GDAL, PROJ), technical flexibility in combining GUI and command-line data processing. The research contributes to the technical cartographic development.}

% Keywords
\keyword{GDAL; PROJ; GRASS GIS; terrain; DEM; elevation; ETOPO1; geomorphometry; cartography; mapping} 

% The fields PACS, MSC, and JEL may be left empty or commented out if not applicable
%\PACS{J0101}
%\MSC{}
%\JEL{}

%%%%%%%%%%%%%%%%%%%%%%%%%%%%%%%%%%%%%%%%%%
\begin{document}

\section{Introduction}

%%%%%%%%%%%%%%%%%%%%%%%%%%%%%%%%%%%%%%%%%%
\section{Materials and Methods}

\subsection{Subsection}

\subsection{Subsection}

\subsection{Subsection}

%%%%%%%%%%%%%%%%%%%%%%%%%%%%%%%%%%%%%%%%%%
\section{Results}

%%%%%%%%%%%%%%%%%%%%%%%%%%%%%%%%%%%%%%%%%%
\section{Discussion}

%%%%%%%%%%%%%%%%%%%%%%%%%%%%%%%%%%%%%%%%%%
\section{Conclusions}

%%%%%%%%%%%%%%%%%%%%%%%%%%%%%%%%%%%%%%%%%%
\authorcontributions{Supervision, conceptualization, methodology, software, resources, funding acquisition, and project administration, O.D.; writing---original draft preparation, methodology, software, data curation, visualization, formal analysis, validation, writing---review and editing, and investigation, P.L. All authors have read and agreed to the published version of the manuscript.}

\funding{The publication was funded by the Editorial Office of XXXX, Multidisciplinary Digital Publishing Institute (MDPI), by providing 100\% discount for the APC of this manuscript. This project was supported by the Federal Public Planning Service Science Policy or Belgian Science Policy Office, Federal Science Policy - BELSPO (B2/202/P2/SEISMOSTORM).}

\institutionalreview{Not applicable.}

\informedconsent{Not applicable.}

\dataavailability{Not applicable} 

\acknowledgments{The authors thank the anonymous reviewers of Land for reading, suggestions and comments that improved an earlier version of this manuscript.}

\conflictsofinterest{The authors declare no conflict of interest.} 

\abbreviations{Abbreviations}{
The following abbreviations are used in this manuscript:\\

\noindent 
\begin{tabular}{@{}ll}
MDPI & Multidisciplinary Digital Publishing Institute\\
DOAJ & Directory of open access journals\\
LD & Linear dichroism
\end{tabular}
}

%%%%%%%%%%%%%%%%%%%%%%%%%%%%%%%%%%%%%%%%%%
%% Optional
\appendixtitles{no} % Leave argument "no" if all appendix headings stay EMPTY (then no dot is printed after "Appendix A"). If the appendix sections contain a heading then change the argument to "yes".
\appendixstart
\appendix
\section[\appendixname~\thesection]{}
\subsection[\appendixname~\thesubsection]{}
The appendix is an optional section that can contain details and data supplemental to the main text---for example, explanations of experimental details that would disrupt the flow of the main text but nonetheless remain crucial to understanding and reproducing the research shown; figures of replicates for experiments of which representative data are shown in the main text can be added here if brief, or as Supplementary Data. Mathematical proofs of results not central to the paper can be added as an appendix.

\begin{table}[H] 
\caption{This is a table caption.\label{tab5}}
\newcolumntype{C}{>{\centering\arraybackslash}X}
\begin{tabularx}{\textwidth}{CCC}
\toprule
\textbf{Title 1}	& \textbf{Title 2}	& \textbf{Title 3}\\
\midrule
Entry 1		& Data			& Data\\
Entry 2		& Data			& Data\\
\bottomrule
\end{tabularx}
\end{table}

\section[\appendixname~\thesection]{}
All appendix sections must be cited in the main text. In the appendices, Figures, Tables, etc. should be labeled, starting with ``A''---e.g., Figure A1, Figure A2, etc.

%%%%%%%%%%%%%%%%%%%%%%%%%%%%%%%%%%%%%%%%%%
\begin{adjustwidth}{-\extralength}{0cm}
%\printendnotes[custom] % Un-comment to print a list of endnotes

\reftitle{References}

%=====================================
% References, variant A: external bibliography
%=====================================
\bibliography{Bibliography.bib}

%%%%%%%%%%%%%%%%%%%%%%%%%%%%%%%%%%%%%%%%%%
\end{adjustwidth}
\end{document}
